\documentclass[11pt]{scrartcl}
\usepackage[parfill]{parskip}
\usepackage{graphics}
\usepackage{hyperref}

\graphicspath{{images/}}

\title{\textbf{Economic Aspects\\
                Exercises}}
\subtitle{Companies running FLOSS business models}
\author{Jesús M. González-Barahona}
\date{\today}

\begin{document}

\maketitle

\vfill

\begin{flushright}
    \copyright  2013 Ricardo Garc\'ia Fern\'andez - ricardogarfe [at] gmail [dot] com.

    This work is licensed under a Creative Commons Attribution- ShareAlike 3.0 License.
    To view a copy of this license visit:
 
    \url{http://creativecommons.org/licenses/by/3.0/legalcode}.
\end{flushright}

\begin{figure}[h]
    \begin{flushright}	
        \includegraphics{by}
        \label{fig:by-sa}
    \end{flushright}
\end{figure}

\newpage

\section{Title}

\section{Description}

Research, using information in the web, the business models of Red Hat, 
Canonical, and Eucalyptus. Describe, briefly, how they earn money, and their 
business strategy. This excercise tryies to be a first approach to business 
models, so no analysis is needed, a brief description is enough.

\subsection{Red Hat}

\quotation{ "The source code is free and it's freely available. But we compile
 those bits and we make it enterprise-class."}\footnote{http://arstechnica.com/business/2012/02/how-red-hat-killed-its-core-productand-became-a-billion-dollar-business/}\\
 
 With this basic principle we can explain Red Hat business model; Enterprise 
 Subscriptions of the product\footnote{http://www.redhat.com/about/subscription/}. 
 This suscription gives to the customer an extra value that doesn't exists
  in dowloadable release, the direct care from RedHat.\\
 The same code, the same freedom but including its extra value. Also their 
 suscription assistance support (standar and premium) they offer courses,
  training and special packaging of every distro (29.9 dolars).\\
 As a conclusion, they profesionalize their FLOSS product making money with 
 these extra options\footnote{http://blip.tv/linuxcom/jim-whitehurst-explains-red-hat-s-business-model-in-less-than-four-minutes-800398}.

\subsection{Canonical}

 Canonical has a closer business model to Red Hat with its distro Ubuntu\footnote{http://www.canonical.com/enterprise-services/ubuntu-advantage/overview}.
 Advantage suscription in Ubuntu provides extra service dedicated to the customer 
 related to specific services included around Ubuntu:
 \quotation{Canonical is to create an ecosystem of products and services around Ubuntu, which would complement the functions of the OS}\footnote{http://www.ghabuntu.com/2009/09/ubuntu-business-model-misunderstood.html}\\
 Without elaborating another side of their business model is partneship to give
 Canonical/Ubuntu support, so selling parnership to other companies and formation
 to these companies is another business model to standarize Ubuntu methodology.\\
 To conclude, we can say that the business model of the use of FLOSS is located
  around the products developed with this technology, support and expertise.
 
\subsection{Eucalyptus}

 Eucalyptus; Open Source Private and Hybrid Clouds from Eucalytus. Here is the 
 slogan from the company, clear on the top of theri website. FLOSS business model
 oriented to the mainstream cloud.
 In section products and services\footnote{http://www.eucalyptus.com/eucalyptus-cloud} 
 we can see what is offer as a product and their partners:
 \begin{itemize}
  \item EUCALYPTUS CLOUD
  \item ECOSYSTEM TOOLS
  \item TRY EUCALYPTUS
  \item EDUCATION
  \item CONSULTING
  \item SUPPORT
\end{itemize}
 Focusing on the support section we can see another sample of FLOSS business model
 as a support service, priority access, specifications and expertise in the product.
 This is very close to other models described above, in another niche but almost the same
 model.\\
 Partnership is another option for the companies that want approvement seal of 
 Eucalyptus in their products, as in other companies (Canonical and Red Hat) they 
 standarize its proccess to expand their knwoledgement and teach it evolving the
  community around their projects creating its work methodology.
 
\section{Conclusion}

As a final conclusion we can say that the business model of each company is 
close to other. Offer extra services to consumers who request them, 
creating standards and partnering with companies that want to export
 the professionalism behind a hallmark of the company itself. 
From what we see as the business model tries through FLOSS free
 base create an ecosystem for the professional who demands 
providing security and reliability by allowing other companies will become 
part of the family of partners. 

\end{document}
